\chapter{\textbf{Input and output}}
\label{chap:importexport}

\section*{Outline}
\addcontentsline{toc}{section}{\texorpdfstring{\hskip14mm}{} Outline}

This chapter covers everything that crosses the boundary of a
\texttt{Cast3M} run: saving and restoring objects, exporting pictures,
reading measurement data, writing results to a file, loading your own
procedures, and -- the section most people arrive here for -- getting a mesh
built in another program into \texttt{Cast3M}.

\section{Saving and restoring \texttt{Cast3M} objects}

\begin{castemcom}
Saving \texttt{Cast3M} objects to an external file is done with
\op{sauv[er]} \dindex{sauv[er]}{save data in a file} (fr. \emph{sauver}).

The option first declares the output file in the working directory. The
string form is what preserves lower and upper case, so the file name comes
out as you wrote it; a directory path may also be given.

\op{sauv} is then followed by the objects to be saved -- here the mesh
\op{plate} and the displacement \op{u}. The program also saves, silently, all
the underlying objects they depend on, which is why the listing reports many
more objects than you named: saving a displacement field drags in its mesh
and all its points.

Called with no argument, \op{sauv;} saves the entire session.
\end{castemcom}
\begin{castemlines}
\begin{verbatim}
opti sauv 'my_file.dat';
sauv plate u;

 LA PILE NUMERO  1 CONTIENT 19 OBJET(S) MAILLAG
  IL Y A  10 OBJET(S) NOMME(S) :
 LA PILE NUMERO  2 CONTIENT   1 OBJET(S) CHPOINT
  IL Y A   1 OBJET(S) NOMME(S) :
 LA PILE NUMERO 27 CONTIENT   1 OBJET(S) MOT
 LA PILE NUMERO 32 CONTIENT 494 OBJET(S) POINT
  IL Y A   6 OBJET(S) NOMME(S) :
 LA PILE NUMERO 33 CONTIENT  1 OBJET(S) CONFIGU
 Fin normale de la sauvegarde
\end{verbatim}
\end{castemlines}

\begin{castemcom}
Data saved with \op{sauv[er]} is recovered with \op{rest[ituer]}
\dindex{rest[ituer]}{recover data from a file}. The option sets the file to
be read and the bare command imports its whole contents; all objects come
back with their previous names.

Check the \texttt{DIMENSION} echoed on the second line of the output: it is
the dimension the objects were saved with, and it must match the
\op{opti dime} of the run reading them.


This pair is the right way to split a long study into stages: mesh in one
file, computation in a second, post-processing in a third. It is also much
faster and completely lossless compared with any exchange format, so for
\texttt{Cast3M}-to-\texttt{Cast3M} transfers always prefer
\op{sauv}/\op{rest} over \op{sort}/\op{lire}.
\end{castemcom}
\begin{castemlines}
\begin{verbatim}
opti rest 'my_file.dat';
rest;

 NIVEAU DU FICHIER 16
 NIVEAU D'ERREUR 0  DIMENSION 2
 DENSITE 0.00000E+00
 LECTURE DE  19 OBJETS MAILLAGE
 LECTURE DE   1 OBJETS CHPOINT
 LECTURE DE   1 OBJETS MOT
 LECTURE DE 494 OBJETS POINT
 LECTURE DE   1 OBJETS CONFIGUR
 FIN DE LECTURE DU LABEL :
 LABEL AUTOMATIQUE :  1
 Fin normale de la restitution
\end{verbatim}
\end{castemlines}

\noindent The default format is binary. \op{opti sauv 'f.dat' 'FORMAT';}
writes a portable text file instead -- slower and larger, but readable and
transferable between machines; \texttt{'XDR'} is the portable binary
alternative. Restoring must use the matching keyword on \op{opti rest}.

\section{Mesh import and export, format exchange}
\label{sec:io:meshexchange}

\cindex{mesh import}\cindex{mesh export}\cindex{gmsh}%
\cindex{MED format}\cindex{UNV format}\cindex{AVS format}%
\cindex{Salome}\cindex{physical group (gmsh)}%
Sooner or later you will want to mesh in a dedicated mesher and compute in
\texttt{Cast3M}. The meshing operators of chapter~\ref{chap:meshing} are
perfectly adequate for the academic geometries of this booklet, but for a
real CAD geometry you want \textsc{gmsh}, \textsc{Salome} or a commercial
pre-processor.

\subsection{What \texttt{Cast3M} can read and write}

Two operators do the work: \op{lire} \dindex{lire}{read an external file} for
import and \op{sort[ir]} \dindex{sort[ir]}{write an external file} for export.
Both act on the file declared by \op{opti lect} and \op{opti sort}
respectively. \textbf{The two lists are not symmetric} -- several formats can
be read but not written, and vice versa:

\begin{center}
\fcolorbox{black}{Gray}{%
\begin{tabular}{@{\hspace{2mm}}>{\ttfamily}l@{\hspace{4mm}}p{62mm}@{\hspace{4mm}}c@{\hspace{3mm}}c@{\hspace{2mm}}}
\multicolumn{4}{c}{\textbf{Mesh and field exchange formats}} \\[2.mm]
 & & \textbf{lire} & \textbf{sort} \\[1.mm]
'MED'  & Salome MED, on HDF5; meshes and fields & $\bullet$ & $\bullet$ \\
'AVS'  & AVS UCD unstructured cell data, ASCII  & $\bullet$ & $\bullet$ \\
'UNV'  & I-DEAS universal file                  & $\bullet$ &           \\
'NAS'  & NASTRAN bulk data                      & $\bullet$ & $\bullet$ \\
'FEM'  & Altair OptiStruct / HyperMesh          & $\bullet$ &           \\
'STL'  & stereolithography, gives a TRI3 mesh   & $\bullet$ & $\bullet$ \\
'CSV'  & plain text columns                     & $\bullet$ &           \\
'ABAQ' & ABAQUS input deck                      &           & $\bullet$ \\
'VTK'  & Paraview \texttt{.vtu} / \texttt{.pvd} &           & $\bullet$ \\
'EXCE' & spreadsheet-readable columns           &           & $\bullet$ \\
\end{tabular}}
\end{center}

\noindent Note what is \emph{not} in that table: there is \textbf{no
\texttt{'GMSH'} keyword}, and \texttt{Cast3M} does not read \textsc{gmsh}
\texttt{.msh} files natively. Note also that \texttt{'UNV'} can be read but
not written -- the FAQ on the website claims otherwise, but the notice of
\op{sort} is authoritative and has no \texttt{'UNV'}. To produce a UNV file,
use the \texttt{cast2unv.pl} script mentioned below.

\subsection{Importing a gmsh mesh}
\label{sec:io:gmsh}

There are three routes. They are listed in the order in which you should try
them.

\paragraph{Route 1 -- via MED (recommended).} MED is the best-supported
exchange format in current \texttt{Cast3M}: the MED library is bundled with
the distribution, it is actively maintained, it is the native format of
\textsc{Salome}, and it carries fields as well as meshes, in both directions.
\textsc{gmsh} writes it directly.

In \textsc{gmsh}, either choose \texttt{File} $\rightarrow$ \texttt{Export}
and give the file a \texttt{.med} extension, or add \texttt{Save "plate.med";}
at the end of your \texttt{.geo} file, or convert an existing mesh from the
command line:

\begin{codebox}
\begin{verbatim}
$ gmsh plate.geo -2 -o plate.med        # mesh and write MED in one go
$ gmsh plate.msh -o plate.med           # convert an existing .msh
\end{verbatim}
\end{codebox}

\begin{castemcom}
\op{lire 'MED'} returns a \op{table}, not a bare mesh: a MED file may hold
several meshes, several element types and several fields, and the table is
indexed by their names.

List it first -- \op{list tab1;} -- to see what the file actually contains,
then pick what you need. This is the step everyone skips and then wonders why
the import ``failed''.

\medskip
Assigning names in the mesher pays off here. In \textsc{gmsh} a
\emph{physical group} named \texttt{TROU} becomes an entry you can address by
that name, which is how you get boundaries you can apply \op{bloq} and
\op{pres} to. Without physical groups you get an anonymous mesh and have to
recover boundaries geometrically with \op{poin} and \op{elem}, as in
chapter~\ref{chap:meshing}.
\end{castemcom}
\begin{castemlines}
\begin{verbatim}
opti dime 2 elem tri3;

tab1 = lire 'MED' 'plate.med';

list tab1;

dom  = tab1 . 'PLAQUE';
trou = tab1 . 'TROU';
bas  = tab1 . 'BAS';

trace (dom et (trou coul rouge));
\end{verbatim}
\end{castemlines}

\paragraph{Route 2 -- via UNV.} If MED is unavailable -- an older
\textsc{gmsh} build without MED support, for instance -- the I-DEAS universal
file is the fallback. \textsc{gmsh} always writes it:

\begin{codebox}
\begin{verbatim}
$ gmsh plate.geo -2 -o plate.unv
\end{verbatim}
\end{codebox}

\noindent and \op{lire 'UNV' 'plate.unv';} reads it. Two cautions. First,
\textsc{gmsh} writes its named groups in UNV dataset 2477, and readers differ
in which group dataset they accept, so check that your named boundaries
actually survived before relying on them. Second, UNV import is read-only --
you cannot write the result back out in the same format.

\paragraph{Route 3 -- the conversion scripts.} The CEA publishes a set of
Perl scripts, in large part due to Laurent Champaney~\cite{Champaney:2007},
on the \textbf{Utilitaires} page of the website, at
\begin{center}
\url{https://www-cast3m.cea.fr/index.php?xml=utilitaires}
\end{center}
\noindent which convert directly:

\begin{center}
\fcolorbox{black}{Gray}{%
\begin{tabular}{@{\hspace{2mm}}>{\ttfamily}l@{\hspace{4mm}}p{78mm}@{\hspace{2mm}}}
\multicolumn{2}{c}{\textbf{Mesh conversion scripts}} \\[2.mm]
msh2gibi.pl & \textsc{gmsh} mesh $\rightarrow$ \texttt{Cast3M} \\
geo2gibi.pl & \textsc{gmsh} geometry $\rightarrow$ \texttt{Cast3M} \\
cast2pos.pl & \texttt{Cast3M} $\leftrightarrow$ \textsc{gmsh} post-processing \\
cast2unv.pl & \texttt{Cast3M} $\rightarrow$ I-DEAS UNV \\
cast2aba.pl & \texttt{Cast3M} $\rightarrow$ ABAQUS \\
bulk2gibi.pl & CATIA/NASTRAN bulk data $\rightarrow$ \texttt{Cast3M} \\
sam2cas.pl  & SAMCEF $\rightarrow$ \texttt{Cast3M} \\
dxf2geo.pl  & DXF $\rightarrow$ \textsc{gmsh} or \texttt{Cast3M} \\
\end{tabular}}
\end{center}

\noindent These produce a \texttt{.dgibi} file you simply include. They are
convenient and they keep group names, but they are community tools of a
certain age -- if one of them chokes on a file, fall back on route 1.

\subsection{A complete worked example}

The following \textsc{gmsh} geometry file reproduces the quarter plate with a
hole of chapter~\ref{chap:meshing}. Put it in \texttt{plate.geo}:

\begin{codebox}
\begin{verbatim}
L = 1.0;  r = 0.1;  lc = 0.03;

Point(1) = {0, 0, 0, lc};         // centre of the hole
Point(2) = {r, 0, 0, lc};
Point(3) = {L, 0, 0, lc};
Point(4) = {L, L, 0, lc};
Point(5) = {0, L, 0, lc};
Point(6) = {0, r, 0, lc};

Line(1) = {2, 3};   Line(2) = {3, 4};
Line(3) = {4, 5};   Line(4) = {5, 6};
Circle(5) = {6, 1, 2};

Curve Loop(1) = {1, 2, 3, 4, 5};
Plane Surface(1) = {1};

// named groups -> named sub-meshes after import
Physical Curve("BAS")    = {1};
Physical Curve("DROITE") = {2};
Physical Curve("HAUT")   = {3};
Physical Curve("GAUCHE") = {4};
Physical Curve("TROU")   = {5};
Physical Surface("PLAQUE") = {1};
\end{verbatim}
\end{codebox}

\noindent Mesh it and export, then read it and carry straight on with the
elastic computation of chapter~\ref{chap:elasticity}:

\begin{codebox}
\begin{verbatim}
$ gmsh plate.geo -2 -o plate.med
\end{verbatim}
\end{codebox}

\begin{castemcom}
Everything downstream is unchanged: once the mesh and its named boundaries
are in \texttt{Cast3M} objects, the model, material, boundary conditions and
solution are exactly the ones of chapter~\ref{chap:elasticity}.

That is the point of this section. Changing the mesher changes the first ten
lines of your file and nothing else.

\medskip
Two things to check on any imported mesh before computing with it: that
\op{opti dime} and \op{opti elem} match what the file contains (a 2D mesh
read with \op{opti dime 3} gives a surface in space), and that the geometry
is where you think it is -- \op{trace} it, and verify the area with
\op{mesu}.
\end{castemcom}
\begin{castemlines}
\begin{verbatim}
opti dime 2 elem tri3;

tab1 = lire 'MED' 'plate.med';
dom  = tab1 . 'PLAQUE';
d12  = tab1 . 'BAS';
d23  = tab1 . 'DROITE';
d34  = tab1 . 'HAUT';
d41  = tab1 . 'GAUCHE';
trou = tab1 . 'TROU';
ldom = 1.;

list (mesu dom 'SURF');

mod = dom mode mecanique
          elastique isotrope;
mat = mod mate 'YOUN' 2.e11 'NU' 0.3;
KK  = rigi mod mat;
\end{verbatim}
\end{castemlines}

\subsection{Exporting a mesh and a field: the AVS round trip}

The simplest exchange format \texttt{Cast3M} both reads and writes is AVS UCD,
a plain-text format. It is worth knowing because it is readable, so it is the
one to reach for when you want to see what was actually exported.

\begin{castemcom}
\op{opti sort[ir]} sets \op{mon\_fichier} as the standard write file, and
\op{sort[ir]} writes the data structure of the mesh \op{dom} -- nodal
coordinates and connectivity -- together with the displacement field \op{u}.
\texttt{'AVS'} is the keyword selecting the output format.
\end{castemcom}
\begin{castemlines}
\begin{verbatim}
opti sort 'mon_fichier';

sort 'AVS' dom u;
\end{verbatim}
\end{castemlines}

\begin{castemcom}
To recover the data, set the reading file with \op{opti lect} and call
\op{lire}, which returns a \op{table}. The indices \op{lemailla} and
\op{lechpoin} hold the mesh (\emph{le maillage}) and the nodal field
(\emph{le champ par points}) respectively.

As always after an import, check the result by plotting both the mesh and one
component of the field rather than assuming the transfer worked.
\end{castemcom}
\begin{castemlines}
\begin{verbatim}
opti lect 'mon_fichier';

tab = lire 'AVS';

trace (tab . lemailla);
trace (tab . lemailla)
      (exco (tab . lechpoin) 'UX');
\end{verbatim}
\end{castemlines}

\subsection{Exporting results to Paraview}

For three-dimensional results, or simply for better pictures than the
built-in viewer gives, export to \textsc{Paraview}:

\begin{castemcom}
\op{sort 'VTK'} writes a \texttt{.vtu} file readable by \textsc{Paraview} and
by most modern visualisation tools. Several time steps can be collected into
a \texttt{.pvd} series.

This is worth knowing even for 2D work: \textsc{Paraview} gives you warped
geometry, arbitrary slices, streamlines, animation and publication-quality
output that the \texttt{Cast3M} window cannot produce.
\end{castemcom}
\begin{castemlines}
\begin{verbatim}
opti sort 'result.vtu';
sort 'VTK' mod u sig;
\end{verbatim}
\end{castemlines}

\subsection{Animating a transient computation}
\label{sec:io:movie}

\cindex{animation}\cindex{time series, exporting}%
A transient computation -- the heating of
chapter~\ref{chap:heat}, or the \op{pasapas} run of
chapter~\ref{chap:plasticity} -- produces one field per time step, and the
natural way to look at it is as an animation. The recipe is to write one
\texttt{.vtu} file per step and collect them into a \texttt{.pvd} series
that \textsc{Paraview} opens as a single time-dependent dataset.

\begin{castemcom}
Inside the time loop, \op{chai[ne]} builds a different file name at each
step and \op{opti sort} redirects to it. The step number must be
zero-padded, or \texttt{res10} will sort before \texttt{res2}.

The \texttt{.pvd} file is a small XML index listing the time value of each
\texttt{.vtu}; writing it from \textsc{gibiane} with \op{mess} is the
simplest route, and the loop below does exactly that.

\medskip
Open \texttt{result.pvd} in \textsc{Paraview}, apply a
\emph{Warp By Vector} filter, press play, and use \emph{Save Animation} to
write a film or an image sequence. For a figure in a report, \emph{Save
Screenshot} at a large resolution is usually better than a film.
\end{castemcom}
\begin{castemlines}
\begin{verbatim}
* -- one .vtu per step --
repeter bsort ndt;
  n   = &bsort;
  nom = chai 'res_' n '.vtu';
  opti sort nom;
  sort 'VTK' mod (U . n) (sig . n);
fin bsort;

* -- the .pvd index --
opti sort 'result.pvd';
mess '<?xml version="1.0"?>';
mess '<VTKFile type="Collection">';
mess '<Collection>';
repeter bpvd ndt;
  n = &bpvd;
  mess (chai '<DataSet timestep="'
             (n * dt)
             '" file="res_' n '.vtu"/>');
fin bpvd;
mess '</Collection>';
mess '</VTKFile>';
opti sort 'sortie.txt';
\end{verbatim}
\end{castemlines}

\noindent The last line redirects the output away again -- forget it and the
rest of your run is written into the \texttt{.pvd} file. The alternative,
and the one to prefer if you are already post-processing in Python, is to
write the plain \texttt{.vtu} files from \texttt{Cast3M} and let
\texttt{meshio} or \texttt{pyvista} assemble the series
(chapter~\ref{chap:python}).

\section{Graphics and plots}
\label{sec:io:graphics}

Graphics created with \op{trac[er]} or \op{dess[iner]} can be saved in
\textsc{postscript} to an external file.
\dindex{trac[er]}{save graphics} \dindex{dess[iner]}{save graphics}
\sindex{save graphics}
The steps are:

\begin{itemize}
\item click, in the \texttt{Cast3M} graphics window, (i) \emph{Softcopy}, then
  (ii) one of \emph{Framemaker}, \emph{Postscript couleur} (colour) or
  \emph{Postscript NB} (black and white);
\item all subsequent graphic output is accumulated in a single temporary file
  \texttt{fort.24} in the working directory. After the end of the run it is
  copied to \texttt{.ps} if the program was started as \texttt{castem26}, or
  to \texttt{my\_file.ps} if it was started as
  \texttt{castem26 my\_file.dgibi}.
\end{itemize}

\noindent The mouse click matters: with this route only the views you actually
display and click get written.

\medskip
In a batch run there is nobody to click, and the second route is the one you
want. Setting the graphics driver with \op{opti 'TRAC'} sends every subsequent
plot straight to the file, with no window and no interaction:

\begin{codebox}
\begin{verbatim}
opti 'TRAC' 'PSC' ;          * colour postscript, no window

trace dom (coor 1 dom);
\end{verbatim}
\end{codebox}

\noindent The output lands in \texttt{.ps}, or in \texttt{my\_file.ps} if the
run was started as \texttt{castem26 my\_file.dgibi}; \op{opti 'FTRA'
'figure.ps';} names it explicitly. The driver keywords are:

\begin{center}
\begin{tabular}{@{}>{\ttfamily}l@{\hspace{6mm}}l@{}}
'X'   & on-screen window (the default) \\
'PS'  & grey-level \textsc{postscript} file \\
'PSC' & colour \textsc{postscript} file \\
'OPEN'& OpenGL window \\
\end{tabular}
\end{center}

\noindent with the full list in the notice of \op{opti}. \op{opti 'TRAC' 'X';}
switches back to the screen.

\subsection*{How to read postscript files}

\textsc{postscript} is a programming language specialised in describing
graphical information, created by \textsc{Adobe}. It is vector-based: the
picture is built from commands rather than from pixels, which is why a
\texttt{Cast3M} figure can be enlarged without degrading, and why its text
stays selectable. Several programs read it:

\begin{itemize}
\item the classical drawing and image-processing packages:
  \textsc{Adobe Illustrator}, \textsc{Adobe Photoshop}, \textsc{Corel Draw},
  \textsc{Inkscape} (free), \textsc{Gimp} (free), \textsc{ImageMagick} (free);
\item the \textsc{ghostscript} family, which is what most viewers use
  underneath. \texttt{ps2pdf} from that family is the quickest way to get a
  figure into a \LaTeX{} document.
\end{itemize}

\noindent For a figure destined for a report, \texttt{ps2pdf figure.ps} and
then \verb|\includegraphics{figure.pdf}| is the whole workflow.

\subsection*{Changing the grey levels in saved graphics}

\textsc{postscript} is a vector language and the files are plain text, so the
grey levels of a saved figure can be edited afterwards. In the header of a
\texttt{Cast3M} figure the definitions look like this:

\begin{castemcom}
The grey levels are defined by their names \op{/eA}, \op{/eB}, \ldots\ and
the associated values \op{1.0000}, \op{0.9841}, \ldots\ respectively.
\op{def} ends the definition. \dindex{trac[er]}{grey levels}

To change the grey scale of a figure, edit these values; the only restriction
is to stay in the $[0,1]$ interval. A global search and replace is usually
enough to turn an unreadable dark figure into a printable one.
\end{castemcom}
\begin{castemlines}
\begin{verbatim}
/eA { 1.0000 setgray } def
/eB { 0.9841 setgray } def
...
/e& { 0.0159 setgray } def
/e@ { 0.0    setgray } def
\end{verbatim}
\end{castemlines}

\section{Reading and writing data}

\subsection{\texttt{acqu[erir]}: reading data from a file}

The standard operator for getting data from a file is \op{acqu[erir]}
\dindex{acqu[erir]}{get data from a file} (fr. \emph{acquérir}, to get). We
illustrate it on a file of measurements.

\begin{castemcom}
The file \texttt{mesures.dat} contains measurements in four columns of real
numbers over 786 lines. Under Linux its structure can be inspected with
\texttt{head}, and \texttt{wc} counts lines, words and characters.
\end{castemcom}
\begin{castemlines}
\begin{verbatim}
$ head mesures.dat
0.    25.58    24.8297   0.7503
0.1   42.4199  24.9874  17.4325
0.2  125.599   24.998  100.601
...
$ wc mesures.dat
    786    3136   27653 mesures.dat
\end{verbatim}
\end{castemlines}

\begin{castemcom}
To read it, the \op{opti acquerir} option is set to the file name, and four
lists are initialised with \op{prog[ramer]}.

The values are then read line by line with \op{acquerir} inside a
\op{repe[ter]} loop. The type of each item is declared with the
\op{*flottant} suffix, and each list grows by appending with \op{et}.

After reading, the lists are assembled into \op{evolution} objects, which can
be displayed with \op{dess[iner]} or interpolated with \op{ipol} -- this is
how a measured material law enters a computation.
\dindex{repe[ter]}{do, start}
\end{castemcom}
\begin{castemlines}
\begin{verbatim}
opti acqu 'mesures.dat';

mes_t    = prog;
mes_text = prog;
mes_tint = prog;
mes_del  = prog;

repeter myloop 786;
  acqu t*flottant text*flottant
       tint*flottant delta*flottant;

  mes_t    = mes_t    et (prog t);
  mes_text = mes_text et (prog text);
  mes_tint = mes_tint et (prog tint);
  mes_del  = mes_del  et (prog delta);
fin myloop;

ev_text  = evol manu 'temps' mes_t
                 't_ext' mes_text;
ev_tint  = evol manu 'temps' mes_t
                 't_int' mes_tint;
ev_delta = evol manu 'temps' mes_t
                 'delta_t' mes_del;

dess (ev_text et ev_tint);
\end{verbatim}
\end{castemlines}

\noindent The alternative, when the file is a clean table, is
\op{lire 'CSV'}, which reads the whole file in one call with options
\texttt{'DEBU'}, \texttt{'FIN'} and \texttt{'SEPA'} for the first line, last
line and separator. It is shorter than the loop above and should be preferred
for new work.

\subsection{The \texttt{@excel1} procedure: writing to a file}

The \op{@excel1} procedure is a user-defined subroutine which saves an
\op{evolution} object in a format readable by a spreadsheet. It was written
in January 1994 by Christian Laborderie (LMT, ENS Cachan) and given to the
\texttt{Cast3M} community.
We analyse it here because it is a good model for any output procedure of
your own -- and because the way it redirects the standard output is the
general mechanism for writing anything to a file from \textsc{gibiane}.

\begin{castemcom}
The complete \textsc{gibiane} structure of a procedure is printed with
\op{list}, which numbers the lines and shows the header \texttt{Cast3M}
stores with it -- name, type, author initials and date.

The procedure is delimited by \op{debp[rocedure]}
\dindex{debp[rocedure]}{subroutine, start} and \op{finproc}. Its input
variables are \op{ev1}, an \op{evolution}, and \op{ficout}, the name of the
output file given as a string (fr. \emph{mot}); the \op{*TYPE} suffix
declares the expected type of each argument.

The two lists of \op{ev1} are extracted with \op{extr[aire]}
\dindex{extr[aire]}{extract from an evolution} and the options
\op{absc[isse]} and \op{ordo[nnee]}.

The current standard output is saved in \op{ii} with \op{vale[ur]}, the
output is redirected to \op{ficout} with \op{opti impr}, the loop prints
pairs of values with \op{mess[age]} \dindex{mess[age]}{print to the standard output}, and the standard
output is finally restored to \op{ii}. Forgetting that last line is the
classic bug: everything the program prints afterwards disappears into your
data file.
\end{castemcom}
\begin{castemlines}
\begin{verbatim}
list @excel1;

DEBPROC @EXCEL1 EV1*EVOLUTION
                FICOUT*MOT;
  PROG1 = EXTR EV1 ABSC;
  PROG2 = EXTR EV1 ORDO;
  NB = DIME PROG1;
  I  = 0;
  ii = VALE IMPR;
  OPTI IMPR 10 IMPR FICOUT;
  REPETER BOU1 NB;
    I  = I + 1;
    X1 = EXTR PROG1 I;
    Y1 = EXTR PROG2 I;
    MESSAGE X1 ';' Y1 ;
  FIN BOU1;
  OPTI IMPR II;
FINPROC;
\end{verbatim}
\end{castemlines}

\noindent This procedure writes a semicolon-separated file of two columns.
It is perfectly serviceable, but a modern workflow exports once and does the
analysis in Python -- which is the subject of chapter~\ref{chap:python}.

\subsection{\texttt{mess[age]} and \texttt{chai[ne]}}

Both of the procedures above print with \op{mess[age]}
\dindex{mess[age]}{print to the standard output}, which writes its arguments
to the current standard output. It accepts numbers and strings in any mixture,
so the quickest way to trace a computation is simply

\begin{codebox}
\begin{verbatim}
mess 'step ' n ' residual ' no_res;
\end{verbatim}
\end{codebox}

\noindent When you need control over the \emph{format} -- a fixed number of
decimals, a value glued into a file name -- convert the number to a string
first with \op{chai[ne]} \dindex{chai[ne]}{convert a value to a string}, which
is also how the CSV writer of chapter~\ref{chap:python} assembles its lines:

\begin{codebox}
\begin{verbatim}
nom = chai 'resultat_' n '.dat';     * build a file name
opti sauv nom;

lig = chai x1 ';' y1;                * build one CSV line
mess lig;
\end{verbatim}
\end{codebox}

\noindent \op{chai} is what lets a loop write a different file at every step,
which is the basis of any parametric study driven from inside
\textsc{gibiane}.

\section{\texttt{util}: loading your own procedures}
\sindex{util[isateur]}

For convenience a set of procedures can be grouped in an external file. Once
that file is loaded, the procedures appear in subsequent runs as if they were
built-in \texttt{Cast3M} commands.

\begin{castemcom}
The file is a plain list of programming lines in which each procedure is
preceded by a line containing \op{\$\$\$\$} followed by the name of the
procedure. The last line of the file contains \op{\$\$\$\$} alone.
\dindex{debp[rocedure]}{procedure file}
\end{castemcom}
\begin{castemlines}
\begin{verbatim}
$$$$ MYFIRST
debproc MYFIRST ...;
*
finproc;

$$$$ MYSECOND
debproc MYSECOND ... ;
finproc;

$$$$
\end{verbatim}
\end{castemlines}

\begin{castemcom}
Historically the file was loaded with the \op{util} directive in a special
run that did nothing else, producing a \texttt{UTILPROC} file in the working
directory that later runs picked up.

\medskip
\textbf{This is no longer the recommended way.} Current versions read
procedures directly from the directories named by the
\texttt{CASTEM\_PROCEDUR26} environment variable, searching the current
directory first, then \texttt{./procedur}, then the installation. The
\op{util} directive is deprecated.

In practice: put your procedures in a \texttt{procedur} subdirectory next to
your \texttt{.dgibi} file, one file per procedure, and they will be found.
Dropping a file there with the name of a standard procedure also overrides
it, which is the supported way to patch one.
\end{castemcom}
\begin{castemlines}
\begin{verbatim}
* -- the old way (deprecated) --
util proc 'my_procedures.dgibi';
fin;

* -- the current way --
*   ./procedur/myfirst.procedur
*   ./procedur/mysecond.procedur
*   and just call them:
myresult = MYFIRST arg1 arg2;
\end{verbatim}
\end{castemlines}

\section{Calling an exterior program}

\op{exte[rieur]} \dindex{exte[rieur]}{call an exterior program} runs an
external program from inside a \texttt{Cast3M} session, and
\op{chau[ssette]} \dindex{chau[ssette]}{dynamic link with another program}
(fr. \emph{chaussette}, socket) opens a socket to another running program,
which is the mechanism used for co-simulation and coupled multiphysics.

The common use of \op{exte} is to hand a file to a script and read the result
back -- for instance to call the Python post-processing of
chapter~\ref{chap:python} at the end of a run, or to drive a parametric study
from outside.

\section*{Summary}
\addcontentsline{toc}{section}{\texorpdfstring{\hskip14mm}{} Summary}

For \texttt{Cast3M}-to-\texttt{Cast3M} transfers use \op{sauv} and
\op{rest}: they are lossless, fast and keep the object names, and they are
what lets you split a study into a meshing file, a computation file and a
post-processing file.

For everything else the two operators are \op{lire} and \op{sort[ir]}, and
their format lists are not symmetric. Coming in from another mesher, MED is
the route to prefer -- it is bidirectional, actively maintained, and it keeps
the group names that become your named boundaries. Going out to a
visualisation tool, \op{sort 'VTK'} and \textsc{Paraview} will do more than
the built-in window ever can. Whatever the route, plot the imported mesh and
check its area before computing on it.

\section*{Exercises}
\addcontentsline{toc}{section}{\texorpdfstring{\hskip14mm}{} Exercises}

\begin{enumerate}
\item \textbf{Round trip.} Mesh the plate with the hole in \textsc{gmsh},
  import it through MED, run the elastic computation, export the result with
  \op{sort 'VTK'} and display the von Mises stress in \textsc{Paraview}.
  Compare the stress concentration with the one obtained on the mesh built
  in chapter~\ref{chap:meshing}.

\item \textbf{Mesh sensitivity across meshers.} Generate the same geometry
  with triangles in \textsc{gmsh} and with quadrangles in \texttt{Cast3M}, at
  comparable numbers of degrees of freedom, and compare the maximum stress.
  Which element type is more accurate here, and why?

\item \textbf{Your own output procedure.} Write a procedure
  \op{@exportcsv} which takes a \op{table} of \op{evolution} objects and
  writes them as a single CSV file with one column per evolution and a header
  line. Use it on the time history of the first exercise of
  chapter~\ref{chap:heat}.
\end{enumerate}
